% GENERATED FILE. Edit canonical lessons, then run npm run generate:books.
% canonical-source-hash: fnv1a64:33040b7640eb5c36
% canonical-lessons: MR-C397-vilamb, MR-C397-kholambne, MR-C397-tatkalne, MR-C397-rakhadne, MR-C397-ravana-hone

\chapter{Held Up on the Way}
\label{ch:mr-held-up-on-the-way}
\hlchaptermodality{397}

\begin{chapteropening}
\textbf{By the end of this chapter:} \emph{I can say a trip was delayed, held up or stalled, and that it set off at last.}

Complete the last lesson of chapter 397: I can say a trip was delayed, held up or stalled, and that it set off at last.
\end{chapteropening}

\section[\texorpdfstring{\mr{विलंब} (vilamb)}{विलंब (vilamb)}]{\mr{विलंब} (vilamb) — a delay}

\label{lesson:MR-C397-vilamb}



\begin{quote}
Before the new one: say the Marathi for a goods train, then the Marathi for to rattle.
\end{quote}

\subsection*{You'll want to know: \mr{विलंब}}
\textbf{\mr{विलंब}} — \emph{vilamb} — \textquotedblleft{}a delay\textquotedblright{}.

\begin{quote}
\mr{पावसामुळे} \mr{गाडीला} \mr{विलंब} \mr{झाला}.
\end{quote}

\begin{quote}
\emph{pāvsāmuḷe gāḍīlā vilamb jhālā.}
\end{quote}

\begin{quote}
The train was delayed because of the rain.
\end{quote}

Masculine.

\mr{विलंब} \mr{होणे}, to be delayed: what is late takes -\mr{ला}, \mr{गाडीला} \mr{विलंब} \mr{झाला}.

\begin{cousinweb}
\textbf{\mr{विलंब}} (\emph{vilamb}) is a tatsama, taken from Sanskrit \emph{vilamba}, \textquotedblleft{}hanging down; slowness, delay\textquotedblright{} (Monier-Williams).
\end{cousinweb}

\begin{grammarlens}[title={What you've built}]
One more word for getting around.
\end{grammarlens}

\subsection*{Your turn}
\begin{itemize}
\raggedright
  \item \emph{Say it:} \emph{vilamb}
  \item \emph{Say it:} \emph{vilamb}, once more
  \item \emph{Say it:} \emph{pāvsāmuḷe gāḍīlā vilamb jhālā.} — The train was delayed because of the rain.
  \item \emph{Recall:} say the Marathi for a goods train, then the Marathi for to rattle, then say \emph{vilamb} again
\end{itemize}

\begin{tcolorbox}[breakable,colback=teal!4,colframe=teal!35!black,title={Before you move on}]
What does \mr{विलंब} mean? (\textquotedblleft{}A delay\textquotedblright{}.) Say it once more.
\end{tcolorbox}
\section[\texorpdfstring{\mr{खोळंबणे} (khoḷambṇe)}{खोळंबणे (khoḷambṇe)}]{\mr{खोळंबणे} (khoḷambṇe) — to be held up}

\label{lesson:MR-C397-kholambne}



\begin{quote}
Before the new one: say the Marathi for to rattle, then the Marathi for a delay.
\end{quote}

\subsection*{You'll want to know: \mr{खोळंबणे}}
\textbf{\mr{खोळंबणे}} — \emph{khoḷambṇe} — \textquotedblleft{}to be held up\textquotedblright{}.

\begin{quote}
\mr{मी} \mr{स्टेशनवर} \mr{दोन} \mr{तास} \mr{खोळंबलो}.
\end{quote}

\begin{quote}
\emph{mī sṭeśanvar don tās khoḷamblo.}
\end{quote}

\begin{quote}
I was held up at the station for two hours.
\end{quote}

Past (he …): \mr{तो} \mr{खोळंबला}, \emph{to khoḷamblā}.

Work can be held up too: \mr{काम} \mr{खोळंबलं}, the work got held up.

\begin{cousinweb}
\textbf{\mr{खोळंबणे}} (\emph{khoḷambṇe}) is a Marathi verb whose origin is not settled; it says that someone or something is stuck waiting.
\end{cousinweb}

\begin{grammarlens}[title={What you've built}]
One more word for getting around.
\end{grammarlens}

\subsection*{Your turn}
\begin{itemize}
\raggedright
  \item \emph{Say it:} \emph{khoḷambṇe}
  \item \emph{Say it:} \emph{khoḷambṇe}, once more
  \item \emph{Say it:} \emph{mī sṭeśanvar don tās khoḷamblo.} — I was held up at the station for two hours.
  \item \emph{Recall:} say the Marathi for to rattle, then the Marathi for a delay, then say \emph{khoḷambṇe} again
\end{itemize}

\begin{tcolorbox}[breakable,colback=teal!4,colframe=teal!35!black,title={Before you move on}]
What does \mr{खोळंबणे} mean? (\textquotedblleft{}To be held up\textquotedblright{}.) Say it once more.
\end{tcolorbox}
\section[\texorpdfstring{\mr{ताटकळणे} (tāṭkaḷṇe)}{ताटकळणे (tāṭkaḷṇe)}]{\mr{ताटकळणे} (tāṭkaḷṇe) — to be kept waiting}

\label{lesson:MR-C397-tatkalne}



\begin{quote}
Before the new one: say the Marathi for a delay, then the Marathi for to be held up.
\end{quote}

\subsection*{You'll want to know: \mr{ताटकळणे}}
\textbf{\mr{ताटकळणे}} — \emph{tāṭkaḷṇe} — \textquotedblleft{}to be kept waiting\textquotedblright{}.

\begin{quote}
\mr{पाहुणे} \mr{दारात} \mr{बराच} \mr{वेळ} \mr{ताटकळले}.
\end{quote}

\begin{quote}
\emph{pāhuṇe dārāt barāc veḷ tāṭkaḷle.}
\end{quote}

\begin{quote}
The guests were kept waiting at the door for a long time.
\end{quote}

Past (he …): \mr{तो} \mr{ताटकळला}, \emph{to tāṭkaḷlā}.

It stresses a long, tiring wait; a woman says \mr{मी} \mr{ताटकळले}.

\begin{cousinweb}
\textbf{\mr{ताटकळणे}} (\emph{tāṭkaḷṇe}) is a Marathi verb for waiting a long time against your will; its origin is not settled.
\end{cousinweb}

\begin{grammarlens}[title={What you've built}]
One more word for getting around.
\end{grammarlens}

\subsection*{Your turn}
\begin{itemize}
\raggedright
  \item \emph{Say it:} \emph{tāṭkaḷṇe}
  \item \emph{Say it:} \emph{tāṭkaḷṇe}, once more
  \item \emph{Say it:} \emph{pāhuṇe dārāt barāc veḷ tāṭkaḷle.} — The guests were kept waiting at the door for a long time.
  \item \emph{Recall:} say the Marathi for a delay, then the Marathi for to be held up, then say \emph{tāṭkaḷṇe} again
\end{itemize}

\begin{tcolorbox}[breakable,colback=teal!4,colframe=teal!35!black,title={Before you move on}]
What does \mr{ताटकळणे} mean? (\textquotedblleft{}To be kept waiting\textquotedblright{}.) Say it once more.
\end{tcolorbox}
\section[\texorpdfstring{\mr{रखडणे} (rakhaḍṇe)}{रखडणे (rakhaḍṇe)}]{\mr{रखडणे} (rakhaḍṇe) — to be stalled}

\label{lesson:MR-C397-rakhadne}



\begin{quote}
Before the new one: say the Marathi for to be held up, then the Marathi for to be kept waiting.
\end{quote}

\subsection*{You'll want to know: \mr{रखडणे}}
\textbf{\mr{रखडणे}} — \emph{rakhaḍṇe} — \textquotedblleft{}to be stalled\textquotedblright{}.

\begin{quote}
\mr{पुलाचं} \mr{काम} \mr{दोन} \mr{वर्षं} \mr{रखडलं}.
\end{quote}

\begin{quote}
\emph{pulācaṁ kām don varṣaṁ rakhaḍlaṁ.}
\end{quote}

\begin{quote}
Work on the bridge was stalled for two years.
\end{quote}

Past (he …): \mr{तो} \mr{रखडला}, \emph{to rakhaḍlā}.

\mr{काम} \mr{रखडणे} is the usual pair: work that stopped and drags on.

\begin{cousinweb}
\textbf{\mr{रखडणे}} (\emph{rakhaḍṇe}) is a Marathi verb for work or plans that stop and drag on; its origin is not settled.
\end{cousinweb}

\begin{grammarlens}[title={What you've built}]
One more word for getting around.
\end{grammarlens}

\subsection*{Your turn}
\begin{itemize}
\raggedright
  \item \emph{Say it:} \emph{rakhaḍṇe}
  \item \emph{Say it:} \emph{rakhaḍṇe}, once more
  \item \emph{Say it:} \emph{pulācaṁ kām don varṣaṁ rakhaḍlaṁ.} — Work on the bridge was stalled for two years.
  \item \emph{Recall:} say the Marathi for to be held up, then the Marathi for to be kept waiting, then say \emph{rakhaḍṇe} again
\end{itemize}

\begin{tcolorbox}[breakable,colback=teal!4,colframe=teal!35!black,title={Before you move on}]
What does \mr{रखडणे} mean? (\textquotedblleft{}To be stalled\textquotedblright{}.) Say it once more.
\end{tcolorbox}
\section[\texorpdfstring{\mr{रवाना} \mr{होणे} (ravānā hoṇe)}{रवाना होणे (ravānā hoṇe)}]{\mr{रवाना} \mr{होणे} (ravānā hoṇe) — to set off; to depart}

\label{lesson:MR-C397-ravana-hone}



\begin{quote}
Before the new one: say the Marathi for to be kept waiting, then the Marathi for to be stalled.
\end{quote}

\subsection*{You'll want to know: \mr{रवाना} \mr{होणे}}
\textbf{\mr{रवाना} \mr{होणे}} — \emph{ravānā hoṇe} — \textquotedblleft{}to set off; to depart\textquotedblright{}.

\begin{quote}
\mr{शेवटी} \mr{बस} \mr{पुण्याकडे} \mr{रवाना} \mr{झाली}.
\end{quote}

\begin{quote}
\emph{śevṭī bas puṇyākaḍe ravānā jhālī.}
\end{quote}

\begin{quote}
At last the bus set off for Pune.
\end{quote}

Past (he …): \mr{तो} \mr{रवाना} \mr{झाला}, \emph{to ravānā jhālā}.

The place you head for takes -\mr{कडे} or -\mr{ला}: \mr{पुण्याला} \mr{रवाना} \mr{होणे}.

\begin{cousinweb}
\textbf{\mr{रवाना} \mr{होणे}} (\emph{ravānā hoṇe}) joins \textbf{\mr{रवाना}}, \textquotedblleft{}sent off, on the way\textquotedblright{}, a loan usually traced to Persian \emph{ravāna}, to the light verb \textbf{\mr{होणे}}, \textquotedblleft{}to become\textquotedblright{}.
\end{cousinweb}

\begin{grammarlens}[title={What you've built}]
One more word for getting around.
\end{grammarlens}

\subsection*{Your turn}
\begin{itemize}
\raggedright
  \item \emph{Say it:} \emph{ravānā hoṇe}
  \item \emph{Say it:} \emph{ravānā hoṇe}, once more
  \item \emph{Say it:} \emph{śevṭī bas puṇyākaḍe ravānā jhālī.} — At last the bus set off for Pune.
  \item \emph{Recall:} say the Marathi for to be kept waiting, then the Marathi for to be stalled, then say \emph{ravānā hoṇe} again
\end{itemize}

\begin{tcolorbox}[breakable,colback=teal!4,colframe=teal!35!black,title={Before you move on}]
What does \mr{रवाना} \mr{होणे} mean? (\textquotedblleft{}To set off; to depart\textquotedblright{}.) Say it once more.
\end{tcolorbox}
